\subsection{Горизонтальный Архимед}
%%%%%%%%%%%%%%%%%%%%%%%%%%%%%%%%%%%%%%%%%%%%%%%%%%%%
%%%%%%%%%%%%%%%%%%%%%%%%%%%%%%%%%%%%%%%%%%%%%%%%%%%%	
\z[{\begin{wrapfigure}[5]{r}{0.2\textwidth}%
		\vspace{-10pt}%
		\begin{flushright}%
 			\includegraphics[width=0.2\textwidth]{9.4.1}%
		\end{flushright}%
	\end{wrapfigure}%
	}]%
	{Тонкая U-образная трубка постоянного внутреннего сечения с горизонтальным коленом длиной $L$ и двумя одинаковыми вертикальными коленами, открытыми в атмосферу, заполнена водой не полностью (см. рисунок). В каждом вертикальном колене остается слой воздуха длиной $H$. Вода начинает выливаться, если трубку двигать вдоль горизонтального колена с постоянным ускорением, не меньшим, чем некоторая величина $a_0$. Горизонтальное колено остаётся всегда заполненным водой. Найти:
	\lpu ускорение $a_0$;\lpuend
	\lpu длину $\Delta L$ вылившегося слоя воды при движении с ускорением $\frac{4a_0}{3}$. \lpuend
	}
%
{\puan  $a_0 = g \dfrac{2H}{L}$;~\puan $x = \dfrac{2}{3}H$}
%%%%%%%%%%%%%%%%%%%%%%%%%%%%%%%%%%%%%%%%%%%%%%%%%%%% 
\z[{\begin{wrapfigure}[7]{r}{0.3\textwidth}%
		\vspace{-10pt}%
		\begin{flushright}%
 			\includegraphics[width=0.3\textwidth]{9.4.2}%
		\end{flushright}%
	\end{wrapfigure}%
	}]%
	{Тонкая запаянная с одного конца трубка заполнена ртутью и закреплена на горизонтальной платформе, вращающейся с угловой скоростью $\omega$ вокруг вертикальной оси так, что ртуть не выливается и заполняет полностью горизонтальное колено трубки (см. рисунок). Открытое колено трубки вертикально. Геометрические размеры установки указаны на рисунке. Атмосферное давление $p_0$, плотность ртути $\rho$.
	\lpu Найти давление $p_1$ ртути в месте изгиба трубки.\lpuend
	\lpu Найти давление $p_2$ ртути у запаянного конца трубки. \lpuend
	}
%
{\puan  $p_1 = p_0 + \rho g H$;~\puan $p_2 = p_0 + \rho g H - \dfrac{3}{2}\rho \omega^2 L^2$.}
%%%%%%%%%%%%%%%%%%%%%%%%%%%%%%%%%%%%%%%%%%%%%%%%%%%%
\z[{\begin{wrapfigure}[5]{r}{0.3\textwidth}%
		\vspace{-15pt}%
		\begin{flushright}%
 			\includegraphics[width=0.3\textwidth]{9.4.3}%
		\end{flushright}%
	\end{wrapfigure}%
	}]%
	{В сосуде с водой находится стеклянный шар. Стенки и дно сосуда гладкие. Дно горизонтальное, левая стенка наклонена под углом $\alpha = 45\degree$ к горизонту (см. рисунок). Объём шара $V$, плотность воды $\rho$, плотность шара $3\rho$.
	\lpu Найдите силу давления $F_1$ шара на дно при неподвижном сосуде.\lpuend
	\lpu Найдите силу давления $F_2$ шара на дно при движении сосуда с горизонтальным ускорением $a = \frac{g}{4}$ \lpuend
	}
%
{\puan  $T_1 = 2\rho g V$;~\puan $T_2 = \dfrac{3}{2}\rho g V$}
%%%%%%%%%%%%%%%%%%%%%%%%%%%%%%%%%%%%%%%%%%%%%%%%%%%%	
\z[{\begin{wrapfigure}[5]{r}{0.3\textwidth}%
		\vspace{-15pt}%
		\begin{flushright}%
 			\includegraphics[width=0.3\textwidth]{9.4.4}%
		\end{flushright}%
	\end{wrapfigure}%
	}]%
	{В сосуде с водой закреплён клин. На гладкой поверхности клина, наклонённой к горизонту под углом $\alpha$ ($\tan \alpha = \frac{1}{4}$), удерживается стеклянный шар с помощью горизонтально натянутой нити (см. рисунок). Объём шара $V$, плотность воды $\rho$, плотность стекла $3\rho$.
	\lpu Найдите силу натяжения $T_1$ нити при неподвижном сосуде. \lpuend
	\lpu Найдите силу натяжения $T_2$ нити при движении сосуда с горизонтальным ускорением $a = \dfrac{g}{8}$. \lpuend

	\noindent
	В обоих случаях шар находится полностью в воде.
	\newpage
	}
%
{\puan  $T_1 = \dfrac{1}{2}\rho g V$; ~\puan $T_2 = \dfrac{3}{4}\rho g V$}
%%%%%%%%%%%%%%%%%%%%%%%%%%%%%%%%%%%%%%%%%%%%%%%%%%%% 
\z[{\begin{wrapfigure}[7]{r}{0.3\textwidth}%
		\vspace{-5pt}%
		\begin{flushright}%
 			\includegraphics[width=0.3\textwidth]{9.4.5}%
		\end{flushright}%
	\end{wrapfigure}%
	}]%
	{На горизонтальной платформе стоит сосуд с водой. В сосуде закреплён тонкий стержень $AB$, наклонённый к горизонту под углом $\alpha$ (рис.). Шар радиусом $R$ может скользить без трения вдоль стержня, проходящего через его центр. Плотность шара $\rho_0$, плотность воды $\rho$ ($\rho_0 < \rho$). При вращении системы с постоянной угловой скоростью вокруг вертикальной оси $OO'$, проходящей через нижний конец $A$ стержня, центр шара устанавливается на расстоянии $l$ от этого конца.
	\lpu С какой силой $N$ шар действует на стержень? \lpuend
	\lpu Найдите угловую скорость $\omega$ вращения платформы. \lpuend
	\lpu При какой минимальной угловой скорости $\omega_{\text{min}}$ вращения шар «утонет» и окажется на дне сосуда? \lpuend
	}
%
{\puan  $N = \dfrac{4}{3}\pi R^3 \dfrac{(\rho - \rho_0)g}{\cos\alpha}$; ~\puan  $\omega = \sqrt{\dfrac{g \sin\alpha}{l \cos^2 \alpha}}$;~\puan $\omega_{\text{мин}} = \sqrt{\dfrac{g}{R}} \tan \alpha$}